\subsection{Sprint 0}

Para essa primeira sprint estavam previstas as seguintes atividades: integração da equipe, visando estabelecer uma base sólida para a colaboração e cooperação necessária entre os integrantes para a continuidade do projeto; a participação ativa da equipe no encontro com os \textit{stakeholders} para a elicitação de requisitos; definição da arquitetura do software, debatendo e selecionando linguagens de programação que mais se enquadram no escopo para serem utilizadas no projeto, além da escolha do banco de dados mais adequado conforme a proposta; preparação da documentação do projeto do framework necessário para sustentar a metodologia definida; desenvolvimento de um protótipo da aplicação, utilizando-se da criação de telas na plataforma Figma com enfoque em design de frontend e na definição da identidade visual do aplicativo; e, por fim, apresentação do protótipo desenvolvido para verificação de alinhamento do projeto com as expectativas dos stakeholders.

Durante esta sprint, a equipe priorizou a realização dessas atividades, cruciais para o sucesso do projeto. Inicialmente, o time empreendeu esforços na confraternização entre os integrantes, visando estabelecer uma base sólida de colaboração fundamentada em um entendimento mútuo e coeso entre os membros e a integração da equipe por meio de um encontro dinâmico que garantiu uma colaboração eficaz e harmoniosa da equipe ao longo do ciclo de desenvolvimento, no qual a aluna participou ativamente.

Juntamente com o time, me preparei com antecedência e participei ativamente no encontro com os \textit{stakeholders}, que garantiu uma clara elicitação de requisitos ao compreender suas necessidades e expectativas e alinhá-las com o escopo. Isso pavimentou o caminho para o desenvolvimento de um projeto eficiente e coeso.

Durante os primeiros encontros presenciais, toda a equipe dedicou-se em debater e, eventualmente, definir cautelosamente a arquitetura do software. Essa etapa fundamental envolveu uma discussão e seleção criteriosa das linguagens de programação que mais se encaixavam com o escopo do projeto e com a metodologia colaborativa adotada pelo time. Além disso, também foi definido o banco de dados mais adequado às necessidades específicas da aplicação. Todo esse processo foi conduzido de forma organizada e eficiente pelo time de AGES IV. 

Os avanços significativos na documentação do projeto e no desenvolvimento do framework também foram notáveis. A abordagem proativa adotada permitiu que o time já estabelecesse uma estrutura sólida preparativa ao desenvolvimento subsequente, otimizando o fluxo de trabalho da equipe e minimizando potenciais obstáculos futuros. 

Ademais, concentrei meus esforços no desenvolvimento de um protótipo de telas na plataforma Figma, trabalhando juntamente com meu time, criei a identidade visual do aplicativo e o design das telas da aplicação na plataforma Figma. Além disso, a equipe realizou a utilização de componentes visando facilitação do processo de desenvolvimento. O time também  concluiu o estabelecimento de um fluxo definido de telas.

Essa abordagem permitiu à equipe alinhar o projeto com as expectativas dos stakeholders. Por fim, a equipe apresentou o protótipo da aplicação ao stakeholder de maneira completa e eloquente, o que contribuiu para a obtenção de um feedback positivo e do alinhamento das ideias e expectativas com a realidade e com o fluxo do projeto.